\begin{theorem}[Ramified prime-cyclic preparation]
\label{mod:ramification-primecyclicpreparation}
Let $K/F$ be a finite valuative extension of nonarchimedean local fields
which is cyclic Galois of prime degree, and suppose that it is ramified:
\[
 e(K/F)=\operatorname{ramificationIndex}(F,K)\ne 1.
\]
The prime-degree dichotomy gives
\[
 f(K/F)=\operatorname{residueDegree}(F,K)=1.
\]

Residue degree one makes the canonical map $k_F\to k_K$ surjective. Every
$F$-automorphism acts trivially on the image of this map, hence trivially on
$k_K$. Since the zeroth lower ramification group is the kernel of the
residue action, this proves the reusable implication
\[
 f(K/F)=1\quad\Longrightarrow\quad G_0=G.
\]

Apply the existing canonical selector
\texttt{PrimeCyclicExtension.lowerBreakIndex} to this equality. Its
boundary theorem supplies a unique $t\in\mathbf N$ satisfying
\[
 \operatorname{IsLowerBreak}_{K/F}(t),\qquad
 G_t=G,\qquad G_{t+1}=1.
\]
Independently, the residue-degree-one monogenic-uniformizer theorem supplies
an integral element $\pi_K\in\mathcal O_K$ with
\[
 \operatorname{IsUniformizer}(\pi_K),\qquad
 \operatorname{Algebra.adjoin}_{\mathcal O_F}\{\pi_K\}
   =\mathcal O_K.
\]
The structure \texttt{PrimeCyclicPreparation} retains exactly the six pieces
of common data
\[
 t,\quad \operatorname{IsLowerBreak}(t),\quad f(K/F)=1,
 \quad \pi_K,\quad \operatorname{IsUniformizer}(\pi_K),
 \quad \mathcal O_F[\pi_K]=\mathcal O_K.
\]
The definition \texttt{primeCyclicPreparation} constructs this structure
from ramifiedness; its selected $t$ is proved equal to the existing canonical
lower-break index.

The tame and wild branches are consequences of the same package. If
$t>0$, the faithful positive-depth ramification coordinate from the norm
filtration preparation gives
\[
 \operatorname{char}(k_F)=[K:F]=e(K/F),
\]
contradicting tame ramification. Thus a tame ramified prime-cyclic
extension has $t=0$, and the package exposes
\[
 \operatorname{IsLowerBreak}_{K/F}(0).
\]

Conversely, at a zero break the faithful degree-zero ramification coordinate
embeds the prime-order group $G_0/G_1$ into $k_K^\times$. Consequently
$[K:F]$ divides $\#k_K-1$, whereas $\operatorname{char}(k_F)$ divides
$\#k_K$. Coprimality of consecutive integers shows that a zero break is
tame. Therefore wild ramification forces $0<t$. The structure
\texttt{WildPrimeCyclicPreparation} extends the common package only by this
positivity proof, and \texttt{wildPrimeCyclicPreparation} constructs it from
ramifiedness and wildness.

The package contains no norm-character finiteness, character conductor,
stationary parameter, Hasse comparison, residual phase, or epsilon-factor
data.
\LeanFile{LanglandsFirstMainLemma/Ramification/PrimeCyclicPreparation.lean}
\lean{LanglandsFirstMainLemma.lowerRamificationGroup_zero_eq_top_of_residueDegree_eq_one}
\lean{LanglandsFirstMainLemma.PrimeCyclicPreparation}
\lean{LanglandsFirstMainLemma.PrimeCyclicPreparation.lowerRamificationGroup_zero_eq_top}
\lean{LanglandsFirstMainLemma.PrimeCyclicPreparation.ramificationIndex_eq_degree}
\lean{LanglandsFirstMainLemma.PrimeCyclicPreparation.t_eq_lowerBreakIndex}
\lean{LanglandsFirstMainLemma.primeCyclicPreparation}
\lean{LanglandsFirstMainLemma.isTamelyRamified_of_isLowerBreak_zero}
\lean{LanglandsFirstMainLemma.PrimeCyclicPreparation.t_eq_zero_of_isTamelyRamified}
\lean{LanglandsFirstMainLemma.PrimeCyclicPreparation.isLowerBreak_zero_of_isTamelyRamified}
\lean{LanglandsFirstMainLemma.PrimeCyclicPreparation.t_pos_of_isWildlyRamified}
\lean{LanglandsFirstMainLemma.WildPrimeCyclicPreparation}
\lean{LanglandsFirstMainLemma.wildPrimeCyclicPreparation}
\leanok
\SourceText{The one-break paragraph preceding equation eq:one-break-new, the tame/wild split in the proof of Theorem thm:norm-filtration, and Section sec:first-main-completion.}
\uses{mod:localfield-monogenicuniformizer,mod:ramification-primecyclicextension,mod:ramification-primedegreedichotomy,mod:ramification-lowergroups,mod:ramification-normbelowbreak}
\FormalizationContract
This node is expected to occupy roughly 400--900 Lean lines. It derives the
residue degree from the live prime-degree dichotomy, proves the depth-zero
lower-group identity from the actual residue action, selects the existing
canonical lower break, and retains the integral uniformizer and its
monogenicity proof as explicit data. The tame and wild specializations are
derived from the canonical break and the live ramification-coordinate API.
No independent copy of a ramification invariant or lower filtration is
introduced.
\end{theorem}
